\chapter{Systemic Infectious Diseases}
\label{ch:infectious}
Infectious diseases remain a leading cause of morbidity and mortality worldwide, disproportionately affecting low- and middle-income countries. Pathogens include bacteria, viruses, parasites, fungi, and prions, each with distinct mechanisms of transmission, pathogenesis, and clinical manifestations. The emergence of antimicrobial resistance, novel pathogens, and climate-driven expansion of vector-borne diseases pose ongoing challenges to global health.

%-------------------------------------------------------------
\section{Bacterial Infections}
%-------------------------------------------------------------

%-------------------------------------------------------------

%-------------------------------------------------------------

%-------------------------------------------------------------

%-------------------------------------------------------------

%-------------------------------------------------------------

%-------------------------------------------------------------

\label{sec:Bacterial Infections}
%-------------------------------------------------------------%-------------------------------------------------------------

\subsection{Anthrax}
\label{sec:Anthrax}
Anthrax is a zoonotic infection caused by \textit{Bacillus anthracis}, a large, gram-positive, spore-forming, encapsulated rod, endemic in agricultural regions worldwide, acquired through contact with infected animals or animal products, and is a category A bioterrorism agent. The condition results from spores entering through skin (cutaneous), inhalation (pulmonary), or ingestion (gastrointestinal), germinating and producing three virulence factors: protective antigen (PA), lethal factor (LF), and edema factor (EF), with LF being a metalloproteinase that cleaves MAPKKs, disrupting signaling and causing macrophage lysis and cytokine storm. Clinical features include cutaneous anthrax (95\% of natural cases): pruritic papule developing into a painless, ulcerated, black eschar surrounded by significant edema (malignant pustule); inhalational anthrax: initial flu-like symptoms followed abruptly by respiratory failure, hemorrhagic mediastinitis, pleural effusions, and shock; gastrointestinal anthrax: abdominal pain, nausea, vomiting, bloody diarrhea, and bowel perforation. Diagnosis is by culture and Gram stain. Management includes anthrax antitoxin plus antibiotics (ciprofloxacin or doxycycline). Prognosis: cutaneous $<$ 1\% mortality; inhalational 40--80\%; gastrointestinal 25--60\%.

\subsection{Botulism (Clostridium botulinum)}
\label{sec:Botulism}
Botulism is a rare, potentially fatal neuroparalytic disease caused by botulinum neurotoxin (BoNT) produced by \textit{Clostridium botulinum}, with approximately 200 cases per year in the US, including foodborne, infant botulism (most common form in US), wound botulism (injection drug users), and iatrogenic forms. The condition results from BoNT (a zinc-dependent metalloproteinase) cleaving SNARE proteins (SNAP-25, synaptobrevin, syntaxin) at the presynaptic neuromuscular junction, preventing acetylcholine release and causing flaccid paralysis. Clinical features include symmetric, descending flaccid paralysis beginning with cranial nerve palsies (diplopia, dysarthria, dysphonia, dysphagia, ptosis) and progressing to descending symmetric weakness of neck, trunk, and extremities, with no sensory deficits and no fever. Diagnosis is clinical, confirmed by detection of BoNT in serum, stool, or food (mouse bioassay or PCR). Management includes equine heptavalent botulism antitoxin (HBAT) for adults or human botulinum immune globulin (BabyBIG) for infants, with supportive care including mechanical ventilation for respiratory failure. Prognosis: mortality $<$ 5\% with modern intensive care.

\subsection{Brucellosis}
\label{sec:Brucellosis}
Brucellosis is a zoonotic infection caused by gram-negative coccobacilli of the genus \textit{Brucella}, endemic in the Mediterranean basin, Middle East, Central Asia, Africa, and Latin America, with 500,000 new cases annually, and transmission through ingestion of unpasteurized dairy products or direct contact with infected animal tissues. The condition results from \textit{Brucella} surviving and replicating within macrophages by inhibiting phagolysosome fusion and producing a type IV secretion system that modulates the host immune response. Clinical features include acute brucellosis (80\%): undulant fever (evening fever spikes that wax and wane), profuse sweating, chills, malaise, arthralgia, myalgia, and back pain, with focal disease including spondylitis, sacroiliitis, peripheral arthritis, epididymo-orchitis, neurobrucellosis, and endocarditis (leading cause of mortality). Diagnosis is by blood culture, bone marrow culture, serology, and PCR. Management includes doxycycline plus rifampin for 6 weeks (standard), with longer therapy for focal disease. Prognosis is excellent for most cases; mortality $<$ 1\% except for endocarditis (up to 20\%).

\subsection{Cholera (Vibrio cholerae)}
\label{sec:Cholera}
Cholera is an acute diarrheal disease caused by \textit{Vibrio cholerae} serogroups O1 and O139, with 1.3--4 million cases and 21,000--143,000 deaths annually worldwide, primarily in developing countries with poor water sanitation. The condition results from colonization of the small intestine and secretion of cholera toxin (CT), an A-B subunit enterotoxin that ADP-ribosylates Gs protein, constitutively activating adenylate cyclase, increasing cAMP, and causing massive chloride and water secretion into the intestinal lumen. Clinical features include sudden onset of profuse, painless, watery diarrhea (rice-water stools), vomiting, and rapidly progressing dehydration leading to hypovolemic shock, metabolic acidosis, and renal failure within hours in severe cases. Diagnosis is based on stool culture (TCBS agar) or rapid dipstick test in endemic settings. Management involves aggressive rehydration with oral rehydration solution (WHO-ORS) for mild to moderate dehydration, intravenous Ringer's lactate for severe dehydration, and antibiotics (doxycycline, azithromycin, or ciprofloxacin) to reduce duration and volume of diarrhea. Prognosis is excellent with prompt rehydration (mortality $<$ 1\%); without treatment, mortality exceeds 50\%.

\subsection{Diphtheria (Corynebacterium diphtheriae)}
\label{sec:Diphtheria}
Diphtheria is a toxin-mediated disease caused by \textit{Corynebacterium diphtheriae}, a gram-positive bacillus, with sporadic cases in high-income countries and outbreaks persisting in regions with low vaccination coverage, transmission by respiratory droplets or contact with infected skin lesions. The condition results from \textit{C. diphtheriae} colonizing the nasopharynx or skin and producing diphtheria toxin, a potent ADP-ribosyltransferase that inactivates elongation factor 2 (EF-2), inhibiting protein synthesis and causing cell necrosis, with toxin causing local tissue destruction (pseudomembrane formation) and systemic effects (myocarditis, neuritis). Clinical features include respiratory diphtheria: sore throat, dysphagia, low-grade fever, and a characteristic grayish-white, thick, adherent pseudomembrane over the tonsils, pharynx, or larynx, with airway obstruction being the primary local risk, and systemic complications including myocarditis and neuropathy. Diagnosis is clinical, with Gram stain and culture confirmation. Management includes diphtheria antitoxin (equine) administered immediately and antibiotics (erythromycin or penicillin G). Prognosis: mortality 5--10\%, lower with antitoxin.

\subsection{Gas Gangrene (Clostridial Myonecrosis)}
\label{sec:Gas Gangrene}
Gas gangrene is a rapidly progressive, life-threatening infection of skeletal muscle caused by \textit{Clostridium perfringens} (80\% of cases) and other clostridial species, with low incidence but high mortality (20--60\%), occurring post-traumatically or spontaneously (associated with colorectal carcinoma, neutropenic colitis). The condition results from spores germinating in devitalized, hypoxic tissue and producing potent exotoxins, including alpha-toxin (lecithinase, phospholipase C) causing hemolysis, platelet aggregation, increased vascular permeability, and direct myonecrosis. Clinical features include sudden onset of severe pain at the wound site, marked edema, tenderness, and crepitus on palpation, with skin progressing from pale to bronze to purplish-black with hemorrhagic bullae, and systemic toxicity including tachycardia, tachypnea, hypotension, fever, and renal failure. Diagnosis is clinical, with Gram stain showing gram-positive rods and radiographs showing gas in soft tissues. Management requires urgent surgical exploration with extensive debridement, high-dose penicillin G plus clindamycin, and aggressive supportive care. Prognosis: mortality is 100\% without treatment; with treatment, 20--60\%.

\subsection{Leprosy (Hansen's Disease)}
\label{sec:Leprosy}
Leprosy is a chronic granulomatous infection caused by \textit{Mycobacterium leprae}, an obligate intracellular, acid-fast bacillus, with approximately 200,000 new cases reported annually, primarily in India, Brazil, and Indonesia, transmitted through prolonged, close contact with untreated patients, and a long incubation period (2--10 years, up to 40 years). The condition results from \textit{M. leprae} invading Schwann cells and macrophages, with the clinical spectrum correlating with the host immune response, ranging from tuberculoid leprosy (strong cell-mediated immunity, few bacilli) to lepromatous leprosy (weak cell-mediated immunity, numerous bacilli). Clinical features include hypopigmented, anesthetic skin patches, peripheral nerve thickening, and neuropathy, with tuberculoid leprosy presenting with 1--3 well-demarcated anesthetic macules/plaques, and lepromatous leprosy presenting with diffuse, symmetric skin infiltration with numerous nodules and plaques (leonine facies), eyebrow loss, and progressive symmetric peripheral neuropathy. Diagnosis is clinical (anesthetic skin patches, nerve enlargement), with skin slit-smear for acid-fast bacilli. Management includes rifampin, dapsone, and clofazimine (multidrug therapy). Prognosis is excellent with treatment.

\subsection{Leptospirosis}
\label{sec:Leptospirosis}
Leptospirosis is a zoonotic spirochetal infection caused by pathogenic serovars of \textit{Leptospira interrogans} sensu lato, with more than 1 million cases annually worldwide, highest incidence in tropical and subtropical regions, and transmission through contact with water or soil contaminated by urine of infected animals. The condition results from \textit{Leptospira} entering through mucous membranes or abraded skin, multiplying in the bloodstream (leptospiremic phase), then localizing to tissues including liver, kidney, meninges, and eye. Clinical features include anicteric leptospirosis (80--90\%): biphasic illness with Phase 1 (leptospiremic: abrupt fever, rigors, severe headache, myalgia, conjunctival suffusion) and Phase 2 (immune: aseptic meningitis, uveitis, rash); Weil's disease (severe icteric leptospirosis, 5--10\%): jaundice, renal failure, hemorrhage, and myocarditis. Diagnosis is by PCR of blood or urine, culture, and serology. Management includes doxycycline or amoxicillin for mild disease, and ceftriaxone or penicillin G for severe disease. Prognosis is excellent for mild disease, but mortality is higher with Weil's disease.

\subsection{Lyme Disease (Borrelia burgdorferi)}
\label{sec:Lyme Disease}
Lyme disease is a tick-borne spirochetal infection caused by \textit{Borrelia burgdorferi} sensu stricto (US) and \textit{B. afzelii} and \textit{B. garinii} (Europe, Asia), and is the most common vector-borne disease in the United States with approximately 500,000 cases annually, transmitted by \textit{Ixodes} ticks requiring a 24--48 hour blood meal for transmission. The condition results from \textit{B. burgdorferi} being injected into the skin with tick saliva, replicating locally, and disseminating hematogenously. Clinical features occur in three stages: Stage 1 (early localized, 3--30 days): erythema migrans---an expanding annular rash with central clearing (``bull's-eye'' lesion) at the tick bite site, often with fever, fatigue, myalgia, and headache; Stage 2 (early disseminated, weeks to months): multiple secondary erythema migrans lesions, flu-like symptoms, migratory arthralgia, cranial nerve palsies, meningitis, and carditis; Stage 3 (late, months to years): Lyme arthritis, acrodermatitis chronica atrophicans, and neuroborreliosis. Diagnosis is by serology using two-tier testing (EIA followed by Western blot). Management includes doxycycline for early disease, ceftriaxone for neurologic disease, and oral antibiotics for Lyme arthritis. Prognosis is excellent with appropriate antibiotic treatment.

\subsection{Necrotizing Fasciitis}
\label{sec:Necrotizing Fasciitis}
Necrotizing fasciitis is a rapidly progressive, life-threatening infection of the subcutaneous tissue and fascia, with secondary involvement of skin and muscle, with 500--1500 cases annually in the US, higher incidence in immunocompromised patients, diabetics, intravenous drug users, and after trauma or surgery, with mortality of 20--40\%. The condition results from infection spreading along the superficial and deep fascia, causing thrombotic occlusion of perforating blood vessels, leading to ischemic necrosis, with Type I (polymicrobial, 70\%) involving anaerobes and facultative organisms, Type II (monomicrobial, 20\%) involving \textit{Streptococcus pyogenes} and \textit{Staphylococcus aureus}, and Type III (gram-negative, marine-related) involving \textit{Vibrio vulnificus}. Clinical features include pain out of proportion to examination (earliest sign), rapidly spreading erythema and edema, induration, skin discoloration, hemorrhagic bullae, crepitus, and systemic toxicity. Diagnosis is clinical, with surgical exploration being paramount. Management requires urgent, aggressive surgical debridement, empirical broad-spectrum antibiotics, and intensive supportive care. Prognosis is guarded, with mortality heavily dependent on prompt surgical intervention.

\subsection{Pertussis (Whooping Cough)}
\label{sec:Pertussis}
Pertussis is a highly contagious respiratory infection caused by \textit{Bordetella pertussis}, a gram-negative coccobacillus, with 24 million cases and 160,000 deaths annually worldwide, with infants $<$ 6 months having the highest morbidity and mortality, and transmission by respiratory droplets. The condition results from \textit{B. pertussis} adhering to ciliated respiratory epithelial cells and producing pertussis toxin, which ADP-ribosylates G proteins and promotes lymphocytosis, with tracheal cytotoxin damaging ciliated epithelial cells. Clinical features include three stages: Catarrhal (1--2 weeks): mild cough, coryza, low-grade fever---most infectious stage; Paroxysmal (2--8 weeks): severe, paroxysmal coughing fits followed by a high-pitched inspiratory whoop, post-tussive emesis, cyanosis, and fatigue; Convalescent (weeks to months): gradual resolution. Diagnosis is by PCR of nasopharyngeal swab (test of choice). Management includes azithromycin (which does not alter the clinical course if begun after the paroxysmal stage but eradicates the organism). Prognosis is excellent with supportive care, though pneumonia is the most common cause of death in infants.

\subsection{Plague (Yersinia pestis)}
\label{sec:Plague}
Plague is a zoonotic infection caused by \textit{Yersinia pestis}, a gram-negative, bipolar-staining coccobacillus, with 1000--5000 cases annually worldwide, endemic in Madagascar, Democratic Republic of Congo, and Peru, transmitted from infected rodents to humans by fleas, and is a category A bioterrorism agent. The condition results from \textit{Y. pestis} being injected by flea bite, phagocytosed by macrophages, and resisting intracellular killing via a type III secretion system, with virulence factors including the F1 capsular antigen and Yops. Clinical features include bubonic plague (most common, 80\%): sudden fever, chills, headache, and tender, swollen, matted lymph nodes (buboes); septicemic plague: high fever, hypotension, DIC, purpuric rash, and acral necrosis; pneumonic plague: fever, productive cough, hemoptysis, and rapidly progressive respiratory failure. Diagnosis is by Gram stain and culture from bubo aspirate, blood, or sputum. Management includes streptomycin or gentamicin, with alternatives including doxycycline, ciprofloxacin, and levofloxacin. Prognosis: bubonic 5--15\% mortality (treated); septicemic 30--50\%; pneumonic nearly 100\% if untreated $>$ 24 hours.

\subsection{Syphilis (Treponema pallidum)}
\label{sec:Syphilis}
Syphilis is a sexually transmitted infection caused by the spirochete \textit{Treponema pallidum} subspecies \textit{pallidum}, with 8 million new cases globally per year, with incidence rising in high-income countries, particularly among men who have sex with men (MSM). The condition results from \textit{T. pallidum} penetrating mucous membranes or abraded skin through sexual contact or transplacentally, with the organism rapidly disseminating hematogenously and lymphatically. Clinical features occur in four stages: Primary (3--90 days): painless chancre (ulcer with clean base, firm, indurated margins) at the site of inoculation, resolving spontaneously in 3--6 weeks; Secondary (2--8 weeks after primary): disseminated rash (non-pruritic, maculopapular, involving palms and soles), condylomata lata, generalized lymphadenopathy, fever, malaise, and mucocutaneous lesions; Latent: asymptomatic; Tertiary (10--30 years): gummatous disease, cardiovascular syphilis (aortitis, aortic regurgitation), and neurosyphilis (general paresis, tabes dorsalis). Congenital syphilis includes stillbirth, rhinitis (``snuffles''), rash, osteochondritis, and Hutchinson's triad. Diagnosis is by dark-field microscopy and serology (non-treponemal tests for screening, treponemal tests for confirmation). Management is benzathine penicillin G (IM) for primary, secondary, and early latent syphilis, with neurosyphilis requiring aqueous crystalline penicillin G IV. Prognosis is excellent with appropriate treatment.

\subsection{Tetanus (Clostridium tetani)}
\label{sec:Tetanus}
Tetanus is a toxin-mediated disease caused by \textit{Clostridium tetani}, a spore-forming, obligate anaerobic bacillus found in soil and animal feces, with approximately 200,000 cases annually worldwide, primarily in low-income countries with inadequate vaccination, and neonatal tetanus accounting for 50\% of tetanus deaths. The condition results from spores germinating in anaerobic wound conditions, producing tetanospasmin (a zinc-dependent metalloproteinase) that binds to presynaptic nerve terminals, is transported retrogradely to the CNS, and cleaves synaptobrevin-2 (VAMP), preventing release of inhibitory neurotransmitters (GABA, glycine) from Renshaw cells, producing unopposed motor neuron activity. Clinical features include an incubation period of 3--21 days, with trismus (lockjaw) being the earliest sign in 75\% of cases, followed by risus sardonicus, opisthotonos, generalized muscle rigidity, and painful spasms triggered by minimal stimuli, with autonomic instability occurring in severe cases. Diagnosis is clinical. Management includes human tetanus immune globulin (HTIG) to neutralize unbound toxin, wound debridement, metronidazole, and supportive care (benzodiazepines for spasms, magnesium sulfate for autonomic instability). Prognosis: mortality remains 10--60\% in severe cases, but prevention is effective through tetanus toxoid vaccination.

\subsection{Tuberculosis (Extrapulmonary)}
\label{sec:Tuberculosis (Extrapulmonary)}
Extrapulmonary tuberculosis accounts for 15--20\% of all tuberculosis cases, occurring when \textit{Mycobacterium tuberculosis} disseminates from the primary pulmonary focus via lymphatic or hematogenous routes, with risk factors including HIV co-infection, immunosuppression, diabetes, and young age. The condition results from primary infection typically occurring in the lungs, with local containment in granulomas, and reactivation or progressive primary disease leading to dissemination to extrapulmonary sites. Clinical features depend on the site: pleuritic involvement (tuberculous pleural effusion, most common EPTB), lymphatic (scrofula---painless, matted, cervical lymphadenopathy), skeletal (Pott's disease, tuberculous arthritis), central nervous system (tuberculous meningitis), genitourinary (sterile pyuria, flank pain), and miliary TB (diffuse hematogenous dissemination with small granulomas throughout the body). Diagnosis is by acid-fast stain and culture of specimens, PCR/Xpert MTB/RIF, and imaging. Management includes standard 6-month regimen (2 months rifampin, isoniazid, pyrazinamide, ethambutol, then 4 months rifampin, isoniazid), with longer therapy for CNS TB and corticosteroids. Prognosis is excellent with appropriate treatment.

\subsection{Typhoid Fever (Salmonella typhi)}
\label{sec:Typhoid Fever}
Typhoid fever is a systemic febrile illness caused by \textit{Salmonella enterica} serotype Typhi, with 11--21 million cases and 200,000 deaths annually worldwide, endemic in South Asia, Southeast Asia, sub-Saharan Africa, and parts of Latin America, with transmission being fecal-oral through contaminated food or water. The condition results from \textit{S.} Typhi invading intestinal M cells, disseminating to mesenteric lymph nodes, and multiplying within macrophages of the reticuloendothelial system before being released into the bloodstream. Clinical features include an incubation of 7--14 days, with gradual onset of fever (stepwise ladder pattern), headache, malaise, dry cough, abdominal pain, relative bradycardia, rose spots, hepatosplenomegaly, and ``pea soup'' diarrhea or constipation, with intestinal hemorrhage or perforation being the classic complication. Diagnosis is by blood culture, bone marrow culture (gold standard), stool culture, and serology. Management includes azithromycin, ceftriaxone, or fluoroquinolones for uncomplicated disease, with MDR and XDR strains increasing. Prognosis is good with appropriate antibiotics; mortality is 10--30\% without treatment.

%-------------------------------------------------------------
\section{Foodborne and Waterborne Parasites}
%-------------------------------------------------------------

%-------------------------------------------------------------

%-------------------------------------------------------------

%-------------------------------------------------------------

%-------------------------------------------------------------

%-------------------------------------------------------------

%-------------------------------------------------------------

\label{sec:Foodborne and Waterborne Parasites}
%-------------------------------------------------------------%-------------------------------------------------------------

\subsection{Cyclospora cayetanensis}
\label{sec:Cyclospora}
Cyclosporiasis is a diarrheal illness caused by \textit{Cyclospora cayetanensis}, a coccidian protozoan parasite, with outbreaks associated with imported fresh produce from endemic regions (Central and South America, the Caribbean, Southeast Asia). The condition results from ingested sporulated oocysts releasing sporozoites in the small intestine that invade enterocytes, causing villous blunting, crypt hyperplasia, and inflammation of the duodenum and jejunum. Clinical features include an incubation of 1--14 days, with prolonged, relapsing watery diarrhea (explosive, non-bloody), nausea, anorexia, weight loss, abdominal cramps, and profound fatigue, with average illness duration of 5--7 weeks if untreated. Diagnosis is by modified Ziehl-Neelsen acid-fast stain of stool or PCR. Management includes trimethoprim-sulfamethoxazole (TMP-SMX). Prognosis is excellent with treatment.

%-------------------------------------------------------------
\section{Fungal Infections (Systemic)}
%-------------------------------------------------------------

%-------------------------------------------------------------

%-------------------------------------------------------------

%-------------------------------------------------------------

%-------------------------------------------------------------

%-------------------------------------------------------------

%-------------------------------------------------------------

\label{sec:Fungal Infections (Systemic)}
%-------------------------------------------------------------%-------------------------------------------------------------

\subsection{Aspergillosis (Invasive)}
\label{sec:Aspergillosis Invasive}
Invasive aspergillosis is a life-threatening infection caused by \textit{Aspergillus} species, most commonly \textit{A. fumigatus}, and is the most common mold infection in immunocompromised patients, with mortality 30--80\%, and risk factors including prolonged neutropenia, HSCT, SOT, high-dose corticosteroids, and advanced HIV/AIDS. The condition results from conidia (spores) being inhaled and reaching the alveoli; in immunocompromised hosts, conidia germinate into hyphae that invade pulmonary blood vessels (angioinvasion), causing thrombosis, tissue infarction, and hemorrhagic necrosis. Clinical features include invasive pulmonary aspergillosis (most common): fever refractory to antibiotics, cough, pleuritic chest pain, and dyspnea, with chest CT showing nodules with halo sign and air crescent sign. Diagnosis is by culture, galactomannan (GM) antigen, and CT chest. Management includes voriconazole (IV or oral) as first-line, with isavuconazole as an alternative. Prognosis is poor without treatment; mortality 30--80\%.

\subsection{Blastomycosis}
\label{sec:Blastomycosis}
Blastomycosis is a systemic fungal infection caused by \textit{Blastomyces dermatitidis}, a dimorphic fungus endemic to the Ohio and Mississippi River Valleys and the Great Lakes region, with 1--2 cases per 100,000 population in endemic areas. The condition results from conidia being inhaled, converting to broad-based budding yeasts in the alveoli, and disseminating hematogenously. Clinical features include pulmonary blastomycosis: acute pneumonia or chronic pneumonia (resembles TB or lung cancer); extrapulmonary blastomycosis: skin is the most common site (verrucous or ulcerative lesions), bone (lytic lesions, osteomyelitis), and genitourinary (prostatitis, epididymitis). Diagnosis is by culture, direct examination of KOH wet mount (large, broad-based budding yeasts), and antigen testing. Management includes itraconazole for mild-moderate pulmonary disease and liposomal amphotericin B for severe, disseminated, or CNS disease. Prognosis is excellent with treatment.

\subsection{Candida auris}
\label{sec:Candida auris}
\textit{Candida auris} is an emerging multidrug-resistant yeast first identified in 2009, causing invasive infections in healthcare settings, with cases reported on six continents, and risk factors including ICU stay, central venous catheters, mechanical ventilation, total parenteral nutrition, recent surgery, and diabetes. The condition results from \textit{C. auris} forming biofilms on medical devices, with mechanisms of antifungal resistance including mutations in ERG11 (azoles), FKS1 (echinocandins), and overexpression of efflux pumps. Clinical features include colonization often being asymptomatic, with invasive disease including candidemia, wound infections, and deep-seated candidiasis. Diagnosis requires MALDI-TOF mass spectrometry or DNA sequencing, as standard biochemical methods misidentify \textit{C. auris}. Management includes echinocandins as first-line, with central line removal essential. Prognosis: mortality 30--60\%.

\subsection{Candidiasis (Invasive / Systemic)}
\label{sec:Candidiasis Systemic}
Invasive candidiasis encompasses candidemia (bloodstream infection) and deep-seated candidiasis (abscesses, endocarditis, endophthalmitis, osteomyelitis, meningitis), and is the most common healthcare-associated fungal infection, with Candida species being the fourth most common cause of bloodstream infections in hospitalized patients. Risk factors include ICU stay, central venous catheters, total parenteral nutrition, broad-spectrum antibiotics, neutropenia, immunosuppression, surgery, and diabetes. The condition results from \textit{Candida} invading through breach of mucocutaneous barriers, adhering to endothelium, and forming biofilms on medical devices. Clinical features include candidemia: fever refractory to broad-spectrum antibiotics, hypotension, and sepsis, often without localizing signs. Diagnosis is by blood culture and (1,3)-$\beta$-D-glucan (BDG) assay. Management includes echinocandins (caspofungin, micafungin, anidulafungin) as first-line, with central line removal being essential. Prognosis is guarded, with significant mortality in ICU patients.

\subsection{Coccidioidomycosis (Valley Fever)}
\label{sec:Coccidioidomycosis}
Coccidioidomycosis (Valley fever) is a systemic fungal infection caused by \textit{Coccidioides immitis} and \textit{C. posadasii}, dimorphic fungi endemic to the Sonoran Desert (southwestern US, northern Mexico, and parts of Central and South America), with 150,000 infections annually in the US, and incidence increasing due to climate change. The condition results from inhalation of arthroconidia leading to conversion into spherules containing endospores in the lung, with a robust Th1 response controlling infection. Clinical features include acute pulmonary: 60\% asymptomatic, 40\% develop Valley fever (influenza-like illness with fever, cough, chest pain, arthralgia, and erythema nodosum); chronic pulmonary: persistent cough, hemoptysis, cavities; disseminated coccidioidomycosis: skin, bones/joints, and CNS (coccidioidal meningitis---the most feared complication). Diagnosis is by culture, serology, and histopathology. Management includes fluconazole for mild-moderate disease and liposomal amphotericin B for severe disseminated disease. Prognosis is excellent for acute pulmonary disease; meningitis requires lifelong therapy.

\subsection{Cryptococcosis}
\label{sec:Cryptococcosis}
Cryptococcosis is a systemic fungal infection caused by \textit{Cryptococcus neoformans} (primarily in immunocompromised) and \textit{Cryptococcus gattii} (emerging, affects immunocompetent individuals), with 250,000 cases annually worldwide and 180,000 deaths, mostly in HIV/AIDS (CD4 $<$ 100), and cryptococcal meningitis being the leading cause of meningitis in HIV-infected adults in sub-Saharan Africa. The condition results from inhalation of basidiospores, with the polysaccharide capsule being the major virulence factor, inhibiting phagocytosis and suppressing T cell responses. Clinical features include cryptococcal meningitis: subacute headache, fever, nausea, and altered mental status, with elevated intracranial pressure being a major cause of morbidity. Diagnosis is by serum cryptococcal antigen (CrAg) and CSF analysis (India ink stain, CrAg). Management includes induction therapy with liposomal amphotericin B plus flucytosine, followed by consolidation and maintenance with fluconazole. Prognosis has improved with modern antifungal therapy but remains poor in advanced HIV/AIDS.

\subsection{Histoplasmosis}
\label{sec:Histoplasmosis}
Histoplasmosis is a systemic fungal infection caused by \textit{Histoplasma capsulatum}, a dimorphic fungus, endemic to the Ohio and Mississippi River Valleys in the United States, and is the most common endemic mycosis in the US, with infection occurring through inhalation of conidia from soil enriched with bat or bird guano. The condition results from conidia being inhaled, reaching the alveoli, and converting to yeast form at body temperature, with yeasts being phagocytosed by macrophages and surviving within phagolysosomes. Clinical features include acute pulmonary histoplasmosis: asymptomatic (90\%) or self-limited flu-like illness; progressive disseminated histoplasmosis (PDH): fever, weight loss, hepatosplenomegaly, lymphadenopathy, pancytopenia, and mucocutaneous ulcers; and chronic pulmonary histoplasmosis: cavitary lung disease resembling TB. Diagnosis is by culture, histopathology, and \textit{Histoplasma} polysaccharide antigen in urine or serum. Management includes itraconazole for mild-moderate disease and liposomal amphotericin B for severe disease. Prognosis is excellent for mild disease; disseminated disease can be fatal without treatment.

\subsection{Mucormycosis}
\label{sec:Mucormycosis}
Mucormycosis is a rare, rapidly progressive, life-threatening fungal infection caused by molds of the order Mucorales (\textit{Rhizopus}, \textit{Mucor}, \textit{Lichtheimia}, \textit{Cunninghamella}), with incidence increasing with the rising number of immunocompromised patients, and risk factors including diabetic ketoacidosis (DKA), neutropenia, HSCT/SOT, iron overload, and corticosteroid use. The condition results from spores being inhaled (rhinocerebral, pulmonary) or inoculated through skin, with the fungus being angioinvasive, causing thrombosis, tissue infarction, and hemorrhage. Clinical features include rhinocerebral mucormycosis (most common, particularly in DKA): fever, facial or periorbital pain, nasal congestion, black necrotic eschar, proptosis, vision loss, and cranial nerve palsies; pulmonary mucormycosis: fever, cough, hemoptysis, and rapid progression. Diagnosis is by tissue biopsy showing broad, ribbon-like, non-septate hyphae with right-angle branching. Management includes urgent surgical debridement and liposomal amphotericin B as first-line antifungal. Prognosis: rhinocerebral 50\% mortality, pulmonary 60\%, disseminated 80--100\%.

%-------------------------------------------------------------
\section{Parasitic Infections}
%-------------------------------------------------------------

%-------------------------------------------------------------

%-------------------------------------------------------------

%-------------------------------------------------------------

%-------------------------------------------------------------

%-------------------------------------------------------------

%-------------------------------------------------------------

\label{sec:Parasitic Infections}
%-------------------------------------------------------------%-------------------------------------------------------------

\subsection{African Trypanosomiasis (Sleeping Sickness)}
\label{sec:African Trypanosomiasis}
African trypanosomiasis (sleeping sickness) is a protozoan infection caused by \textit{Trypanosoma brucei gambiense} (West African, 98\% of cases, chronic) and \textit{T. b. rhodesiense} (East African, acute), transmitted by the tsetse fly, with fewer than 2000 cases annually, endemic in sub-Saharan Africa. The condition results from trypomastigotes being inoculated by tsetse fly bite, multiplying in the blood, lymph, and interstitial spaces (hemolymphatic stage), with the parasite undergoing antigenic variation by switching its variant surface glycoprotein (VSG) coat, eventually crossing the blood-brain barrier. Clinical features include hemolymphatic stage (Stage 1): chancre at the tsetse fly bite site, intermittent fever, headache, joint pain, generalized lymphadenopathy (Winterbottom's sign), and pruritus; meningoencephalitic stage (Stage 2): progressive confusion, personality changes, daytime somnolence, nighttime insomnia, headache, tremor, ataxia, and finally coma and death. Diagnosis is by microscopic identification of motile trypomastigotes in blood, lymph node aspirate, or CSF. Management depends on staging: Stage 1 is treated with pentamidine or suramin; Stage 2 is treated with nifurtimox-eflornithine combination therapy (NECT) or melarsoprol. Prognosis is good with early treatment; Stage 2 disease has higher mortality.

\subsection{Amebiasis (Entamoeba histolytica)}
\label{sec:Amebiasis}
Amebiasis is an intestinal and extraintestinal infection caused by \textit{Entamoeba histolytica}, a protozoan parasite, with 50 million symptomatic cases and 100,000 deaths annually worldwide, particularly in the Indian subcontinent, sub-Saharan Africa, and Central/South America. The condition results from cysts excysting in the terminal ileum, releasing trophozoites that colonize the large intestine, with \textit{E. histolytica} trophozoites adhering to colonic mucin, secreting cysteine proteinases that degrade extracellular matrix, inducing host cell apoptosis, and invading the colonic epithelium. Clinical features include intestinal amebiasis: amoebic dysentery---bloody mucoid diarrhea (strawberry jelly stools), abdominal pain, tenesmus, and weight loss; extraintestinal amebiasis: amebic liver abscess---most common extraintestinal form with fever, right upper quadrant pain, and ``anchovy paste'' pus. Diagnosis is by fresh stool microscopy (trophozoites with ingested RBCs), antigen detection, PCR, and serology. Management includes tinidazole or metronidazole followed by a luminal agent (paromomycin) to eradicate cysts. Prognosis is excellent with treatment.

\subsection{Chagas Disease (American Trypanosomiasis)}
\label{sec:Chagas Disease}
Chagas disease is a protozoan infection caused by \textit{Trypanosoma cruzi}, transmitted by reduviid (triatomine) bugs (``kissing bugs''), with 6--8 million infected, primarily in Latin America, with increasing cases in the US, Europe, and Japan due to migration. The condition results from metacyclic trypomastigotes entering the skin or mucous membranes, invading macrophages and other nucleated cells, and transforming into amastigotes, with the parasite having a tropism for cardiac muscle, smooth muscle, and autonomic ganglia. Clinical features include acute phase (6--8 weeks, often asymptomatic): Romãna's sign (unilateral, painless palpebral edema and conjunctivitis), fever, lymphadenopathy, hepatosplenomegaly, and occasionally acute myocarditis; indeterminate phase: asymptomatic, positive serology; chronic phase (10--30 years after infection, 30\% progress): Chagas cardiomyopathy---the most common and serious manifestation, megaesophagus, and megacolon. Diagnosis is by microscopic identification of trypomastigotes in blood (acute) or serology (chronic). Management includes benznidazole or nifurtimox for acute and congenital disease, with supportive care for chronic cardiomyopathy. Prognosis is good for acute disease; chronic cardiomyopathy has significant morbidity.

\subsection{Filariasis (Lymphatic Filariasis)}
\label{sec:Filariasis}
Lymphatic filariasis (elephantiasis) is a mosquito-borne nematode infection caused by \textit{Wuchereria bancrofti} (90\% of cases), \textit{Brugia malayi}, and \textit{B. timori}, with 120 million infected, primarily in tropical and subtropical regions, and is the second leading cause of permanent disability worldwide. The condition results from third-stage filariform larvae (L3) being inoculated by mosquito bite, migrating to the lymphatics, and maturing into adult worms in lymphatic vessels and lymph nodes, with the host immune response to dying adult worms causing granulomatous inflammation and lymphatic damage. Clinical features include asymptomatic microfilaremia (most infected individuals), acute filarial fevers (adenolymphangitis with painful lymphadenitis, fever, chills), and chronic manifestations (lymphedema of the leg, scrotal hydrocele, elephantiasis). Diagnosis is by microscopic detection of microfilariae on thick blood smear or antigen detection. Management includes diethylcarbamazine (DEC) and albendazole. Prognosis is good with early treatment; chronic lymphedema requires long-term care.

\subsection{Giardiasis}
\label{sec:Giardiasis}
Giardiasis is a protozoan diarrheal illness caused by the flagellate \textit{Giardia lamblia}, the most common intestinal parasitic infection in the United States and a leading cause of waterborne outbreaks globally, with 300 million cases annually worldwide. The condition results from cysts excysting in the duodenum, releasing trophozoites that adhere to enterocytes via a ventral sucking disk, causing enterocyte damage, microvillus atrophy, and increased intestinal permeability. Clinical features include an incubation of 1--3 weeks, with acute giardiasis presenting with watery, foul-smelling diarrhea (sulfuriferous), abdominal cramps, bloating, flatulence, and steatorrhea, and chronic giardiasis presenting with persistent diarrhea, weight loss, malabsorption, and fatigue. Diagnosis is by microscopy of stool for cysts and trophozoites or antigen detection. Management includes tinidazole or metronidazole. Prognosis is excellent with treatment.

\subsection{Hookworm / Soil-Transmitted Helminths}
\label{sec:Hookworm Soil-Transmitted Helminths}
Soil-transmitted helminths include hookworm (\textit{Ancylostoma duodenale}, \textit{Necator americanus}), roundworm (\textit{Ascaris lumbricoides}), and whipworm (\textit{Trichuris trichiura}), constituting the most common infections of humans globally, with 1.5 billion people infected, endemic in tropical and subtropical regions with poor sanitation. The condition results from hookworm larvae penetrating skin, migrating through the circulation to the lungs, ascending the trachea, and maturing in the small intestine, with adult worms attaching to intestinal mucosa and feeding on blood. Clinical features include hookworm: cutaneous larva currens, iron deficiency anemia (from chronic blood loss), fatigue, pallor, and growth retardation; \textit{Ascaris}: heavy infection causes abdominal pain, malnutrition, and intestinal obstruction; \textit{Trichuris}: dysentery syndrome (bloody, mucoid stools, tenesmus, rectal prolapse). Diagnosis is by stool microscopy with Kato-Katz method. Management includes albendazole or mebendazole. Prognosis is excellent with treatment.

\subsection{Leishmaniasis}
\label{sec:Leishmaniasis}
Leishmaniasis is a protozoan infection caused by \textit{Leishmania} species, transmitted by sandflies, with a clinical spectrum ranging from localized cutaneous disease to lethal visceral leishmaniasis (Kala-azar), with 1.5 million new cases annually and 20,000--30,000 deaths, endemic in tropical and subtropical regions. The condition results from the sandfly injecting promastigotes into the skin, which are phagocytosed by macrophages and convert to amastigotes, replicating intracellularly and spreading to the reticuloendothelial system. Clinical features include visceral leishmaniasis: fever (often double quotidian), weight loss, massive splenomegaly, hepatomegaly, lymphadenopathy, and pancytopenia; cutaneous leishmaniasis: single or multiple painless ulcer(s) at sandfly bite sites with raised, indurated borders; mucocutaneous leishmaniasis: destructive lesions of the nasopharynx. Diagnosis is by microscopic identification of amastigotes in tissue aspirates or biopsies. Management includes liposomal amphotericin B (first-line for visceral), pentavalent antimonials, and miltefosine. Prognosis depends on the clinical form and treatment.

\subsection{Malaria (Plasmodium spp.)}
\label{sec:Malaria}
Malaria is a protozoan infection caused by \textit{Plasmodium} species, transmitted by female \textit{Anopheles} mosquitoes, with five species causing human disease: \textit{P. falciparum} (most virulent, responsible for nearly all deaths), \textit{P. vivax}, \textit{P. ovale}, \textit{P. malariae}, and \textit{P. knowlesi}, with 240 million cases and 600,000 deaths annually, primarily in sub-Saharan Africa. The condition results from sporozoites being injected with mosquito saliva, traveling to the liver and invading hepatocytes (exoerythrocytic schizogony), with merozoites released from the liver invading RBCs and initiating the erythrocytic cycle, with \textit{P. falciparum} modifying RBC surfaces and causing cytoadherence, microvascular obstruction, and end-organ damage. Clinical features include an incubation of 7--30 days, with paroxysms of fever, chills, rigor, and diaphoresis every 48--72 hours depending on species, with uncomplicated malaria presenting with cyclic fever, headache, myalgia, anemia, and splenomegaly, and severe malaria (from \textit{P. falciparum}) presenting with cerebral malaria, severe anemia, acute kidney injury, pulmonary edema, hypoglycemia, metabolic acidosis, DIC, and shock. Diagnosis is by thick and thin blood smears and rapid diagnostic tests. Management includes artemisinin-based combination therapy (ACT) for uncomplicated \textit{P. falciparum}, artesunate IV for severe malaria, and chloroquine plus primaquine for \textit{P. vivax} and \textit{P. ovale}. Prognosis is excellent with prompt treatment; severe malaria has significant mortality.

\subsection{Primary Amebic Meningoencephalitis (Naegleria fowleri)}
\label{sec:Naegleria fowleri}
Primary amebic meningoencephalitis is a rare but nearly always fatal CNS infection caused by \textit{Naegleria fowleri}, a free-living amoeba found in warm freshwater environments, with fewer than 20 cases annually in the United States, primarily in children and young adults with a history of swimming in warm freshwater lakes. The condition results from the trophozoite form invading the olfactory neuroepithelium, migrating along the olfactory nerve, and entering the subarachnoid space, producing a hemorrhagic, necrotizing meningoencephalitis. Clinical features include an incubation of 1--9 days, with fulminant onset of severe frontal headache, fever, nausea, vomiting, and meningismus, with rapid progression to altered mental status, seizures, coma, and death within 3--7 days. Diagnosis is by PCR of CSF. Management includes miltefosine plus fluconazole, azithromycin, and rifampin, with therapeutic hypothermia and management of increased intracranial pressure. Prognosis: survival is extremely rare (fewer than 15 survivors documented worldwide).

\subsection{Schistosomiasis}
\label{sec:Schistosomiasis}
Schistosomiasis (bilharziasis) is a trematode infection caused by blood flukes of the genus \textit{Schistosoma}: \textit{S. haematobium} (urinary), \textit{S. mansoni} (intestinal), \textit{S. japonicum}, \textit{S. intercalatum}, and \textit{S. mekongi} (intestinal), with 200 million infected and 800 million at risk, primarily in sub-Saharan Africa, with transmission through contact with fresh water containing infected snails. The condition results from cercariae penetrating human skin, transforming into schistosomula, migrating to the portal venous system, maturing into adult worms, and producing eggs that become trapped in tissues, causing granulomatous inflammation, fibrosis, organomegaly, and obstruction. Clinical features include acute schistosomiasis (Katayama fever): fever, cough, urticaria, myalgia, and hepatosplenomegaly 2--6 weeks after exposure; chronic \textit{S. haematobium}: hematuria, dysuria, hydronephrosis, and bladder fibrosis, with squamous cell carcinoma as a long-term complication; chronic \textit{S. mansoni} and \textit{S. japonicum}: hepatosplenomegaly, periportal fibrosis (Symmers' pipe-stem fibrosis), portal hypertension, and esophageal varices. Diagnosis is by microscopic identification of eggs in stool or urine. Management includes praziquantel. Prognosis is good with treatment.

\subsection{Tapeworm Infections (Taeniasis / Cysticercosis)}
\label{sec:Tapeworm Infections}
Taeniasis is intestinal infection with adult tapeworms (\textit{Taenia saginata}---beef tapeworm; \textit{Taenia solium}---pork tapeworm; \textit{Diphyllobothrium latum}---fish tapeworm), and cysticercosis is tissue infection with the larval stage of \textit{T. solium}, a leading cause of acquired epilepsy in endemic areas, with taeniasis affecting 50 million people and neurocysticercosis affecting 2.5 million. The condition results from taeniasis: ingestion of raw or undercooked meat containing cysticerci, with the scolex attaching to the intestinal mucosa; cysticercosis: ingestion of \textit{T. solium} eggs, with oncospheres penetrating the intestinal wall, disseminating hematogenously, and encysting in tissues. Clinical features include taeniasis: usually asymptomatic or mild abdominal pain and passage of proglottids; cysticercosis: neurocysticercosis---seizures (most common presentation), headache, focal neurologic deficits, increased intracranial pressure, and hydrocephalus. Diagnosis of taeniasis is by macroscopic or microscopic identification of proglottids in stool; diagnosis of cysticercosis is by neuroimaging showing ring-enhancing cysts. Management includes praziquantel or niclosamide for taeniasis, and albendazole plus corticosteroids for neurocysticercosis. Prognosis is excellent for taeniasis; neurocysticercosis prognosis depends on cyst burden and location.

\subsection{Toxoplasmosis}
\label{sec:Toxoplasmosis}
Toxoplasmosis is a protozoan infection caused by \textit{Toxoplasma gondii}, an obligate intracellular parasite, with approximately 30\% of the global population seropositive, with cats being the definitive host, and humans acquiring infection through ingestion of oocysts or tissue cysts, or transplacentally. The condition results from oocysts or tissue cysts releasing bradyzoites/tachyzoites in the GI tract, with tachyzoites disseminating hematogenously and infecting all nucleated cells, with the host cellular immune response driving tachyzoites into dormant bradyzoite forms within tissue cysts. Clinical features in immunocompetent hosts are usually asymptomatic or self-limited mononucleosis-like illness; immunocompromised hosts (HIV/AIDS) develop toxoplasmic encephalitis---most common cerebral mass lesion in AIDS---presenting with headache, confusion, seizures, and ring-enhancing lesions on MRI; congenital toxoplasmosis presents with the classic triad of chorioretinitis, hydrocephalus, and intracranial calcifications. Diagnosis is by serology, PCR, and MRI for CNS disease. Management includes pyrimethamine plus sulfadiazine plus folinic acid for active disease. Prognosis is excellent for immunocompetent hosts; congenital infection has significant morbidity.

%-------------------------------------------------------------
\section{Prion Diseases}
%-------------------------------------------------------------

%-------------------------------------------------------------

%-------------------------------------------------------------

Prion diseases, or transmissible spongiform encephalopathies, are rare, fatal neurodegenerative conditions caused by misfolded prion proteins ($PrP^{Sc}$). Misfolded proteins induce conformational changes in native cellular prion proteins ($PrP^C$), driving neurotoxic aggregation and vacuolar encephalopathy. Clinical features include rapidly progressive dementia, myoclonus, ataxia, and severe neurological decline, carrying a universally fatal prognosis.

%-------------------------------------------------------------

%-------------------------------------------------------------

%-------------------------------------------------------------

%-------------------------------------------------------------

\label{sec:Prion Diseases}
%-------------------------------------------------------------%-------------------------------------------------------------

\subsection{Creutzfeldt-Jakob Disease (CJD)}
\label{sec:Creutzfeldt-Jakob Disease}
Creutzfeldt-Jakob disease is a rapidly progressive, fatal neurodegenerative disease caused by prions---misfolded, protease-resistant isoforms (PrP$^{Sc}$) of the normal cellular prion protein (PrP$^{C}$), with an incidence of 1--2 per million annually worldwide, median age at onset 60--65 years, with sporadic CJD (sCJD) accounting for 85--90\% of cases, familial CJD (10\%) linked to PRNP mutations, and iatrogenic CJD ($<$1\%) from contaminated surgical instruments. The condition results from PrP$^{Sc}$ converting native PrP$^{C}$ into the misfolded conformation through a template-directed process, leading to exponential accumulation of PrP$^{Sc}$ aggregates, causing vacuolation (spongiform change), neuronal loss, and astrocytosis. Clinical features include rapidly progressive dementia (weeks to months), myoclonus (startle-sensitive, characteristic), cerebellar ataxia, pyramidal and extrapyramidal signs, visual disturbances, and akinetic mutism in the terminal stage, with death typically occurring within 1 year of symptom onset. Diagnosis is by clinical criteria, EEG showing periodic sharp wave complexes (PSWCs), CSF RT-QuIC (specific and sensitive), and MRI showing T2/FLAIR hyperintensity in the basal ganglia and thalamus. Management is palliative. Prognosis is uniformly fatal.

\subsection{Kuru}
\label{sec:Kuru}
Kuru is a prion disease historically endemic among the Fore people of Papua New Guinea, transmitted through ritualistic endocannibalism, with peak incidence in the 1950s--1960s, and the last known case reported in 2009. The disease was the first human prion disease shown to be transmissible (to chimpanzees, by Carleton Gajdusek, 1966). The condition results from the same mechanism as other prion diseases, with the incubation period being inversely related to the dose of infectious material and age at exposure. Clinical features include three stages: ambulant (tremor and progressive ataxia, ``kuru'' means ``to shiver'' in Fore), sedentary (severe ataxia, dysarthria, loss of independent ambulation), and terminal (complete motor incapacity, primitive reflexes, incontinence, and mutism), with emotional lability and inappropriate euphoria (``kuru laughter'') being characteristic. Diagnosis is clinical (historical). Management is palliative. Prognosis: the disease was eliminated by the cessation of endocannibalism.

\subsection{Variant CJD (Mad Cow Disease)}
\label{sec:Variant CJD}
Variant Creutzfeldt-Jakob disease is a prion disease caused by transmission of bovine spongiform encephalopathy (BSE, ``mad cow disease'') to humans through ingestion of BSE-contaminated beef products, with 232 cases reported worldwide, mean age at onset 28 years (much younger than sCJD), and a longer duration (14 months). The condition results from the same mechanism as sCJD (PrP$^{Sc}$ replication and neurotoxicity) but with a distinct prion strain and unique glycoform banding pattern, with PrP$^{Sc}$ also found in lymphoid tissue, unlike sCJD. Clinical features include early psychiatric symptoms (depression, anxiety, withdrawal, hallucinations) predominating for several months, followed by cerebellar ataxia, cognitive decline, myoclonus, choreoathetosis, and akinetic mutism. Diagnosis is by MRI showing the pulvinar sign (bilateral high signal in the posterior thalamus) and tonsil biopsy detecting PrP$^{Sc}$. Management is palliative. Prognosis is uniformly fatal.

%-------------------------------------------------------------
\section{Sepsis and Septic Shock}
%-------------------------------------------------------------

%-------------------------------------------------------------

%-------------------------------------------------------------

%-------------------------------------------------------------

%-------------------------------------------------------------

%-------------------------------------------------------------

%-------------------------------------------------------------

\label{sec:section:Sepsis and Septic Shock}
%-------------------------------------------------------------%-------------------------------------------------------------

\subsection{Sepsis and Septic Shock}
\label{sec:Sepsis and Septic Shock}
Sepsis is life-threatening organ dysfunction caused by a dysregulated host response to infection, with septic shock being a subset with profound circulatory, cellular, and metabolic abnormalities associated with a higher risk of mortality, with an estimated 48.9 million cases and 11 million deaths annually globally. The condition results from infection triggering pathogen-associated molecular patterns (PAMPs) and damage-associated molecular patterns (DAMPs) that bind toll-like receptors (TLRs), leading to release of pro-inflammatory cytokines (TNF-$\alpha$, IL-1, IL-6), neutrophil activation, complement cascade, endothelial injury, and microvascular thrombosis, with mitochondrial dysfunction and immune exhaustion contributing to late mortality. Clinical features include fever or hypothermia, tachycardia, tachypnea, altered mental status, hypotension, oliguria, mottled skin, and evidence of end-organ dysfunction (elevated lactate, acute kidney injury, hepatic dysfunction, thrombocytopenia). Diagnosis is by clinical criteria per Sepsis-3 guidelines (suspected infection plus SOFA score increase of $\ge$ 2 points), with blood cultures, lactate, complete blood count, coagulation profile, renal and hepatic function, procalcitonin, and imaging to identify the source. Management includes early recognition, source control, empirical broad-spectrum antibiotics (within 1 hour of recognition), intravenous fluids (30 mL/kg crystalloids for hypotension), vasopressors (norepinephrine first-line), and supportive care. Prognosis depends on the underlying infection, timeliness of treatment, and patient factors.

%-------------------------------------------------------------
\section{Sexually Transmitted Infections}
%-------------------------------------------------------------

%-------------------------------------------------------------

%-------------------------------------------------------------

%-------------------------------------------------------------

%-------------------------------------------------------------

%-------------------------------------------------------------

%-------------------------------------------------------------

\label{sec:Sexually Transmitted Infections}
%-------------------------------------------------------------%-------------------------------------------------------------

\subsection{Chancroid}
\label{sec:Chancroid}
Chancroid is a sexually transmitted infection caused by \textit{Haemophilus ducreyi}, a fastidious gram-negative coccobacillus, most common in resource-limited settings, and a cofactor for HIV transmission due to genital ulceration. The condition results from \textit{H. ducreyi} entering through microabrasions during sexual contact and producing a cytolethal distending toxin that causes epithelial cell necrosis. Clinical features include an incubation of 4--10 days, with painful, non-indurated genital ulcers with ragged, undermined borders and a purulent, necrotic base, and painful inguinal lymphadenopathy occurring in 50\% of cases. Diagnosis is clinical, with Gram stain showing gram-negative coccobacilli in ``school of fish'' arrangement. Management includes azithromycin, ceftriaxone, ciprofloxacin, or erythromycin. Prognosis is excellent with treatment.

\subsection{Lymphogranuloma Venereum (LGV)}
\label{sec:Lymphogranuloma Venereum}
Lymphogranuloma venereum is a sexually transmitted infection caused by \textit{Chlamydia trachomatis} serovars L1, L2, and L3, which are invasive and disseminate to regional lymph nodes, endemic in Africa, Southeast Asia, and the Caribbean, with outbreaks among MSM in Europe, North America, and Australia. The condition results from LGV serovars infecting epithelial cells and macrophages, then disseminating through lymphatics, causing lymphangitis and perilymphangitis with granulomatous inflammation. Clinical features include an incubation of 3--30 days, with primary stage: painless papule or ulcer at the inoculation site; secondary stage (2--6 weeks): painful inguinal lymphadenopathy with the ``groove sign'' characteristic; tertiary stage: genital elephantiasis and rectal strictures. Diagnosis is by \textit{C. trachomatis} PCR from ulcer, rectal swab, or bubo aspirate, with genotyping distinguishing LGV serovars. Management includes doxycycline 100 mg twice daily for 21 days. Prognosis is excellent with treatment.

\subsection{Non-Gonococcal Urethritis (NGU)}
\label{sec:Non-Gonococcal Urethritis}
Non-gonococcal urethritis is urethral inflammation without \textit{Neisseria gonorrhoeae} infection, the most common urethral syndrome in sexually active men, with \textit{Chlamydia trachomatis} (serovars D--K) accounting for 30--50\% of cases, \textit{Mycoplasma genitalium} (15--25\%), \textit{Trichomonas vaginalis} (5--10\%), and other causes. The condition results from infection of the urethral epithelium triggering a neutrophilic inflammatory response. Clinical features include an incubation of 2--35 days, with dysuria, urethral discharge (mucoid to mucopurulent), and urethral pruritus, with many infections being asymptomatic, and complications including epididymitis, prostatitis, and reactive arthritis. Diagnosis is by first-void urine NAAT for \textit{C. trachomatis}, \textit{N. gonorrhoeae}, \textit{M. genitalium}, and \textit{T. vaginalis}. Management includes azithromycin or doxycycline for chlamydia, with moxifloxacin as second-line for \textit{M. genitalium}. Prognosis is excellent with treatment.

%-------------------------------------------------------------
\section{Tickborne and Vector-Borne Diseases}
%-------------------------------------------------------------

%-------------------------------------------------------------

%-------------------------------------------------------------

%-------------------------------------------------------------

%-------------------------------------------------------------

%-------------------------------------------------------------

%-------------------------------------------------------------

\label{sec:Tickborne and Vector-Borne Diseases}
%-------------------------------------------------------------%-------------------------------------------------------------

\subsection{Eastern Equine Encephalitis}
\label{sec:Eastern Equine Encephalitis}
Eastern equine encephalitis is a rare but severe mosquito-borne viral disease caused by Eastern equine encephalitis virus (EEEV), an alphavirus, and is the most severe arboviral encephalitis in North America, with 5--10 human cases annually in the United States and a case fatality rate of 30--50\%. The condition results from EEEV crossing the blood-brain barrier and infecting neurons, causing neuronal necrosis, perivascular cuffing, and intense neutrophilic inflammation. Clinical features include an incubation of 4--10 days, with a systemic phase (sudden fever, chills, headache, myalgia, malaise) followed by an encephalitic phase (altered mental status, seizures, cranial nerve palsies, and focal deficits), with rapid progression to coma and respiratory failure. Diagnosis is by EEEV-specific IgM in CSF or serum by ELISA. Management is supportive care. Prognosis: 30--50\% mortality; survivors have a high rate (50--80\%) of severe neurologic sequelae.

\subsection{Ehrlichiosis}
\label{sec:Ehrlichiosis}
Ehrlichiosis is a tick-borne illness caused by obligate intracellular bacteria of the family Anaplasmataceae, with human monocytic ehrlichiosis (HME) caused by \textit{Ehrlichia chaffeensis}, transmitted by the Lone Star tick (\textit{Amblyomma americanum}), with over 1500 cases annually in the United States. The condition results from \textit{E. chaffeensis} infecting monocytes and macrophages, forming intracytoplasmic microcolonies (morulae), and inhibiting phagolysosomal fusion. Clinical features include an incubation of 5--14 days, with fever, severe headache, myalgia, malaise, nausea, and rigors, with rash being less common (30\%) than in RMSF, and laboratory findings of leukopenia, thrombocytopenia, and elevated hepatic transaminases being characteristic. Diagnosis is by PCR of blood (high sensitivity in the first week). Management includes doxycycline. Prognosis is excellent with treatment; mortality 1--3\%.

\subsection{Rocky Mountain Spotted Fever}
\label{sec:Rocky Mountain Spotted Fever}
Rocky Mountain spotted fever is a tick-borne illness caused by \textit{Rickettsia rickettsii}, an obligate intracellular gram-negative coccobacillus, and is the most severe tick-borne rickettsial disease in the Americas, with 2000--4000 cases annually in the United States. The condition results from \textit{R. rickettsii} infecting vascular endothelial cells, causing widespread endothelial damage through direct cytotoxicity and immune-mediated injury, leading to increased vascular permeability, edema, coagulation activation, and microvascular thrombosis. Clinical features include an incubation of 2--14 days, with the classic triad of fever, severe headache, and rash appearing on days 2--5, beginning as blanching erythematous macules on the wrists, ankles, and forearms, spreading centripetally to the trunk and palms/soles, with severe complications including encephalitis, acute kidney injury, ARDS, and gangrene. Diagnosis is clinical in endemic areas---treatment should not be delayed for laboratory confirmation. Management includes doxycycline. Prognosis: mortality 20--30\% if untreated; 3--5\% with appropriate treatment.

%-------------------------------------------------------------
\section{Viral Infections}
%-------------------------------------------------------------

%-------------------------------------------------------------

%-------------------------------------------------------------

Viral infections encompass a broad spectrum of human pathogens capable of obligate intracellular replication within host tissues. Viral entry, genomic integration, and cellular lysis or immune-mediated destruction trigger systemic inflammatory responses and tissue injury. Therapeutic strategies rely on preventive vaccination, antiviral medications targeting viral replication enzymes, and supportive care.

%-------------------------------------------------------------

%-------------------------------------------------------------

%-------------------------------------------------------------

%-------------------------------------------------------------

\label{sec:Viral Infections}
%-------------------------------------------------------------%-------------------------------------------------------------

\subsection{Arboviral Encephalitis}
\label{sec:Arboviral Encephalitis}
Arboviral encephalitis encompasses viral infections transmitted by arthropod vectors that cause brain inflammation, with major causes in North America including La Crosse encephalitis, St. Louis encephalitis, and Western equine encephalitis.

\textbf{La Crosse Encephalitis:} Caused by La Crosse virus (bunyavirus), transmitted by \textit{Aedes triseriatus}. It is the most common cause of pediatric arboviral encephalitis in the United States (50--150 cases annually). Clinical features: incubation 5--15 days, with most infections being asymptomatic or mild febrile illness; neuroinvasive disease presents with fever, headache, seizures, altered mental status, and focal deficits. Case fatality rate is less than 1\%.

\textbf{St. Louis Encephalitis:} Caused by St. Louis encephalitis virus (flavivirus), transmitted by \textit{Culex} mosquitoes. Older adults have the highest risk of neuroinvasive disease. Clinical features: incubation 4--21 days, with neuroinvasive disease presenting with fever, headache, nuchal rigidity, altered mental status, tremor, cranial nerve palsies, and hyponatremia. Case fatality rate is 5--15\%.

\textbf{Western Equine Encephalitis:} Caused by Western equine encephalitis virus (alphavirus), transmitted by \textit{Culex} mosquitoes. Infants are most severely affected. Case fatality rate is 3--5\% in children and up to 15\% in adults.

Diagnosis of arboviral encephalitis is by serology (IgM in CSF or serum by ELISA). Management is supportive care. No specific antiviral therapy is available.

\subsection{Chickenpox}
\label{sec:Chickenpox}
Chickenpox (varicella) is a highly contagious primary infection caused by varicella-zoster virus (VZV, human herpesvirus 3), predominantly affecting unvaccinated children. The condition results from airborne respiratory droplet transmission or direct contact with vesicular fluid, with viral replication in nasopharyngeal lymph nodes followed by viremic dissemination to skin and viscera. Clinical features include a brief viral prodrome (fever, malaise, pharyngitis) followed by a characteristic intensely pruritic maculopapular rash that rapidly evolves into clear vesicles on an erythematous base ('dewdrops on a rose petal'), which rupture and crust over 48 hours. Lesions appear in successive crops over 3--5 days, resulting in lesions at all stages of development simultaneously across the body. Complications include secondary bacterial skin infections (Staphylococcus aureus, Streptococcus pyogenes), cerebellar ataxia, pneumonia (more severe in adults), and encephalitis; VZV remains dormant in dorsal root ganglia and may reactivate as herpes zoster later in life. Diagnosis is clinical, confirmed by Tzanck smear (multinucleated giant cells) or PCR of lesion fluid. Management in healthy children is supportive (antipruritics, cool compresses; aspirin is strictly contraindicated due to Reye syndrome); oral acyclovir or valacyclovir is indicated for adolescents, adults, and immunocompromised patients. Live attenuated varicella vaccine provides strong protection.

\subsection{Chikungunya}
\label{sec:Chikungunya}
Chikungunya is a mosquito-borne viral illness caused by chikungunya virus (CHIKV), an alphavirus transmitted by \textit{Aedes aegypti} and \textit{A. albopictus}, with outbreaks in Africa, Asia, the Indian Ocean islands, Europe, and the Americas. The condition results from CHIKV infecting fibroblasts, epithelial cells, and monocytes, activating immune cells that produce cytokines mediating acute fever and severe arthralgia, with the virus persisting in macrophages and synovial tissue, contributing to chronic arthropathy. Clinical features include an incubation of 2--6 days, with abrupt onset of high fever, severe polyarthralgia (symmetric, involving hands, wrists, ankles, knees, and spine), maculopapular rash on trunk and limbs, and myalgia, with chronic arthralgia ($>$ 3 months) occurring in 30--60\% of patients. Diagnosis is by RT-PCR (first 5 days of illness) and serology. Management is supportive (rest, hydration, acetaminophen for fever, NSAIDs for arthralgia after dengue is excluded). Prognosis is good for acute disease; chronic arthralgia can be debilitating.

\subsection{Cytomegalovirus (CMV)}
\label{sec:Cytomegalovirus}
Cytomegalovirus is a beta-herpesvirus that infects 60--90\% of the adult population worldwide, generally asymptomatic in immunocompetent hosts but causes significant disease in immunocompromised patients and congenitally infected neonates, and is the most common congenital infection (0.5--1\% of live births). The condition results from CMV infecting epithelial cells, endothelial cells, fibroblasts, and monocytes, with primary infection followed by lifelong latency and reactivation occurring with immunosuppression. Clinical features depend on the host: immunocompetent host presents with mononucleosis-like syndrome; immunocompromised host (HIV with CD4 $<$ 50, transplant recipients) presents with CMV retinitis, CMV pneumonitis, CMV colitis, and CMV esophagitis; congenital CMV presents with sensorineural hearing loss (most common sequela), microcephaly, intracranial calcifications, chorioretinitis, and hepatosplenomegaly. Diagnosis is by PCR (quantitative) of blood, urine, CSF, or BAL fluid. Management includes ganciclovir or valganciclovir for immunocompromised patients and congenital CMV. Prognosis depends on the host immune status.

\subsection{Dengue Fever}
\label{sec:Dengue Fever}
Dengue fever is a mosquito-borne viral illness caused by dengue virus (DENV 1--4), a flavivirus transmitted by \textit{Aedes aegypti} mosquitoes, with 400 million infections and 100 million symptomatic cases annually, and 40,000 deaths, endemic in tropical and subtropical regions. The condition results from the virus infecting dendritic cells, keratinocytes, and macrophages, with a secondary infection with a different serotype being the primary risk factor for severe dengue (antibody-dependent enhancement, ADE). Clinical features include incubation of 3--14 days, with classic dengue fever presenting with sudden onset high fever, severe headache (retro-orbital), arthralgia and myalgia (``breakbone fever''), maculopapular rash, lymphadenopathy, and leukopenia; severe dengue (dengue hemorrhagic fever/dengue shock syndrome) presents with plasma leakage, severe thrombocytopenia, mucosal bleeding, and shock during the critical phase at defervescence. Diagnosis is by NS1 antigen test, IgM and IgG serology, and RT-PCR. Management is supportive (acetaminophen for fever, avoiding NSAIDs and aspirin due to bleeding risk, fluid resuscitation). Prognosis is excellent for classic dengue; severe dengue has significant mortality without prompt supportive care.

\subsection{Ebola Virus Disease}
\label{sec:Ebola Virus Disease}
Ebola virus disease is a severe, often fatal hemorrhagic fever caused by Ebola virus, with five species: Zaire (most virulent), Sudan, Bundibugyo, Taï Forest, and Reston, with the 2014--2016 West Africa outbreak being the largest (28,646 cases, 11,323 deaths), and transmission through direct contact with blood, secretions, or bodily fluids of infected humans or animals. The condition results from the virus entering through mucosal surfaces or breaks in skin, infecting dendritic cells, macrophages, and monocytes, spreading to lymph nodes, liver, spleen, and adrenal glands, causing profound immunosuppression and release of pro-inflammatory cytokines (cytokine storm), with endothelial cell dysfunction, DIC, and multi-organ failure. Clinical features include an incubation of 2--21 days, with prodrome (sudden fever, severe headache, myalgia, malaise), gastrointestinal phase (severe watery diarrhea, nausea, vomiting, abdominal pain), and systemic phase (multi-organ failure, coagulopathy, hypotension/shock, and neurologic involvement). Diagnosis is by RT-PCR of blood. Management includes supportive care and two FDA-approved monoclonal antibody therapies: mAb114 (ansuvimab) and REGN-EB3. Prognosis: Zaire species 50--80\% mortality without treatment; 33--50\% with supportive care and monoclonal antibodies.

\subsection{Epstein-Barr Virus / Infectious Mononucleosis (Systemic Aspects)}
\label{sec:Epstein-Barr Virus Systemic}
Epstein-Barr virus is a gamma-herpesvirus that infects $>$ 90\% of the global population, with primary infection causing infectious mononucleosis (IM) in adolescents and young adults, and EBV being associated with several malignancies: Burkitt lymphoma, Hodgkin lymphoma, post-transplant lymphoproliferative disorder, nasopharyngeal carcinoma, and gastric carcinoma. The condition results from EBV infecting B lymphocytes via CD21 (CR2) receptor, leading to polyclonal B cell proliferation, with the host T cell response producing the atypical lymphocytosis characteristic of IM. Clinical features include an incubation of 30--50 days, with the classical triad of fever, pharyngitis (severe, exudative), and cervical lymphadenopathy, along with profound fatigue, splenomegaly, hepatomegaly, and rash (especially with ampicillin/amoxicillin administration). Systemic complications include airway obstruction, splenic rupture, hemolytic anemia, thrombocytopenia, myocarditis, and encephalitis. Diagnosis is by atypical lymphocytosis, Monospot test, and EBV-specific antibodies. Management is supportive. Prognosis is excellent for acute IM.

\subsection{Hand, Foot, and Mouth Disease}
\label{sec:Hand, Foot, and Mouth Disease}
Hand, foot, and mouth disease (HFMD) is a highly contagious pediatric viral exanthem caused by enteroviruses, most commonly Coxsackievirus A16 and Enterovirus 71. The condition results from fecal-oral or oral-oral transmission with primary viral replication in pharyngeal and gastrointestinal lymphoid tissue followed by viremic cutaneous and mucosal dissemination. Clinical features include low-grade fever, malaise, and sore throat followed by painful vesiculoulcerative enanthem on the buccal mucosa, tongue, and soft palate, and a non-pruritic, non-tender maculopapular or oval vesicular exanthem on the palms of the hands, soles of the feet, and buttocks. Diagnosis is clinical based on characteristic lesion distribution. Management is supportive with oral fluid hydration and analgesics (aspirin is strictly avoided); severe neurological complications (aseptic meningitis, brainstem encephalitis) may rarely occur with Enterovirus 71.

\subsection{Hantavirus Pulmonary Syndrome}
\label{sec:Hantavirus Pulmonary Syndrome}
Hantavirus pulmonary syndrome is a severe, often fatal respiratory illness caused by hantaviruses, with the primary cause in North America being Sin Nombre virus (SNV), with 800 cases reported in the Americas since 1993, mortality 30--40\%, and transmission through inhalation of aerosolized excreta from infected rodents. The condition results from the virus infecting endothelial cells (particularly in the lungs), causing increased vascular permeability without direct cytotoxicity, mediated by CD8+ T cell activation and cytokine release disrupting endothelial barrier function. Clinical features include an incubation of 9--35 days, with prodrome (3--5 days): fever, myalgia, headache, and gastrointestinal symptoms, followed by rapid onset of cardiopulmonary phase: cough, dyspnea, and respiratory failure from pulmonary edema, often within 24 hours, with hypotension and shock being common. Diagnosis is clinical in endemic areas, confirmed by serology. Management is aggressive supportive care (mechanical ventilation, vasopressors, ECMO in severe cases). Prognosis: mortality 30--40\%.

\subsection{Human Metapneumovirus (hMPV)}
\label{sec:Human Metapneumovirus}
Human metapneumovirus is a single-stranded RNA virus (family Paramyxoviridae) first identified in 2001, causing acute respiratory tract infections across all ages, with the highest burden in young children, older adults, and immunocompromised patients, accounting for 5--15\% of hospitalized respiratory illnesses in children. The condition results from hMPV infecting ciliated epithelial cells, causing syncytia formation, ciliary dysfunction, and epithelial sloughing, with the virus inhibiting type I interferon signaling. Clinical features include an incubation of 3--6 days, with presentation nearly indistinguishable from RSV: infants and young children present with coryza, cough, fever, wheezing, bronchiolitis, and pneumonia; older adults and immunocompromised present with severe lower respiratory tract infection. Diagnosis is by RT-PCR of nasopharyngeal swab or BAL. Management is supportive care. Prognosis is excellent for immunocompetent hosts.

\subsection{Marburg Virus}
\label{sec:Marburg Virus}
Marburg virus disease is a severe hemorrhagic fever caused by Marburg virus, a negative-sense RNA virus (genus \textit{Marburgvirus}), the only filovirus other than Ebola to cause human disease, with outbreaks in Africa, and the reservoir being the Egyptian fruit bat (\textit{Rousettus aegyptiacus}). The condition results from identical pathophysiology to Ebola virus---infection of dendritic cells, macrophages, and monocytes, profound immunosuppression, cytokine storm, endothelial dysfunction, and DIC. Clinical features include an incubation of 2--21 days, with abrupt onset of fever, chills, severe headache, myalgia, and prostration, followed by gastrointestinal phase (severe watery diarrhea, abdominal pain) and hemorrhagic phase (petechiae, ecchymoses, mucosal bleeding, uncontrolled hemorrhage, and multi-organ failure). Diagnosis is by RT-PCR (blood, urine, throat swab). Management is supportive care with no approved specific therapy. Prognosis: mortality 23--80\% depending on outbreak and strain.

\subsection{Measles (Rubeola)}
\label{sec:Measles}
Measles is a highly contagious vaccine-preventable viral illness caused by measles virus, a paramyxovirus, with 140,000 deaths globally in 2018, and one of the most contagious infectious diseases with a basic reproduction number ($R_0$) of 12--18. The condition results from the virus entering through respiratory epithelium, replicating in regional lymph nodes, and causing viremia, spreading to the reticuloendothelial system, with infection causing profound lymphopenia and immunosuppression lasting weeks to months. Clinical features include an incubation of 7--14 days, with prodrome (2--4 days): fever, the three Cs---cough, coryza, conjunctivitis, and Koplik spots (pathognomonic: small, white, bluish-white papules on the buccal mucosa), followed by rash (day 3--5): erythematous, maculopapular, confluent rash beginning on the face and spreading cephalocaudally. Complications include pneumonia (most common cause of death), otitis media, laryngotracheobronchitis, encephalitis, and subacute sclerosing panencephalitis (SSPE). Diagnosis is clinical, confirmed by RT-PCR or IgM serology. Management is supportive with vitamin A. Prognosis is excellent with supportive care; prevention through MMR vaccine is highly effective.

\subsection{Mpox}
\label{sec:Mpox}
Mpox is a viral zoonosis caused by mpox virus, a double-stranded DNA orthopoxvirus closely related to variola virus, formerly called monkeypox, endemic in Central and West Africa, with the 2022 global outbreak involving over 90,000 cases across 110 countries. The condition results from the virus entering through broken skin, respiratory tract, or mucous membranes, replicating at the inoculation site, spreading to regional lymph nodes, and causing viremia. Clinical features include an incubation of 5--21 days, with prodrome (fever, headache, myalgia, lymphadenopathy---hallmark distinguishing mpox from smallpox) followed by rash appearing 1--4 days after fever, progressing from macules to papules, vesicles, pustules, and crusts over 2--4 weeks, with a centrifugal distribution, and anogenital lesions common in the 2022 outbreak. Diagnosis is by PCR of lesion swabs. Management includes tecovirimat (TPOXX) and supportive care. Prognosis is generally good; clade I (Central Africa) is more virulent than clade II (West Africa).

\subsection{Mumps}
\label{sec:Mumps}
Mumps is an acute viral illness caused by mumps virus, a paramyxovirus, with 500,000 cases annually worldwide before widespread vaccination, now less common in vaccinated populations but with outbreaks in close-contact settings. The condition results from the virus entering through respiratory epithelium, replicating in the upper respiratory tract and regional lymph nodes, and causing viremia, with tropism for glandular tissue and the CNS. Clinical features include an incubation of 16--18 days, with prodrome (fever, headache, myalgia, malaise) followed by parotitis (bilateral in 70--90\%): painful, tender swelling of the parotid glands, with complications including orchitis (post-pubertal males, 20--30\%), oophoritis, meningitis, pancreatitis, and sensorineural hearing loss. Diagnosis is clinical, confirmed by RT-PCR of buccal swab or IgM serology. Management is supportive. Prognosis is excellent; prevention through MMR vaccine is highly effective.

\subsection{Norovirus}
\label{sec:Norovirus}
Norovirus is a single-stranded RNA virus (family Caliciviridae) and the leading cause of acute gastroenteritis worldwide, responsible for 685 million cases and 200,000 deaths annually, with all age groups affected and peaks in winter months (``winter vomiting disease''). The condition results from the virus infecting enterocytes of the small intestine, causing blunting and atrophy of intestinal villi, crypt hyperplasia, and increased apoptosis, reducing absorptive capacity. Clinical features include an incubation of 12--48 hours, with sudden onset of nausea, vomiting (more prominent in children), watery diarrhea (non-bloody), abdominal cramps, and low-grade fever, with symptoms typically lasting 24--72 hours. Diagnosis is by RT-PCR of stool. Management is supportive care with oral rehydration therapy. Prognosis is excellent in immunocompetent hosts.

\subsection{Poliomyelitis}
\label{sec:Poliomyelitis}
Poliomyelitis is a highly infectious viral disease caused by poliovirus (serotypes 1, 2, 3) of the \textit{Enterovirus} genus, with endemic transmission persisting in Afghanistan and Pakistan. The condition results from fecal-oral transmission, viral replication in the oropharynx and gut, and hematogenous spread to the central nervous system, where it infects and destroys anterior horn cells of the spinal cord and motor nuclei of the brainstem. Clinical features include asymptomatic infection ($>$ 90\%), abortive poliomyelitis (febrile illness with sore throat, nausea), non-paralytic aseptic meningitis (1--5\%), and paralytic poliomyelitis (0.1--2\%) presenting with asymmetric, flaccid paralysis of lower limbs more than upper limbs, areflexia, without sensory loss. Bulbar paralysis affects cranial nerves and respiratory muscles. Diagnosis is based on stool viral culture, RT-PCR of stool or CSF, and serology. Management is supportive: physical therapy, orthotics, respiratory support for bulbar involvement, and prevention of complications. Prognosis for paralytic poliomyelitis: 10\% die from respiratory paralysis, 30\% recover some function, 60\% have permanent paralysis. Prevention is through OPV (Sabin) or IPV (Salk) vaccination; global eradication is 99.9\% complete.

\subsection{Rabies}
\label{sec:Rabies}
Rabies is a fatal viral encephalitis caused by rabies virus, a single-stranded RNA virus of the genus \textit{Lyssavirus}, with 60,000 deaths annually worldwide, primarily in Asia and Africa, with dogs being the main reservoir and source of transmission (99\% of human cases), and rabies being nearly 100\% fatal after symptom onset. The condition results from the virus being inoculated through saliva of an infected animal, binding to nicotinic acetylcholine receptors at the neuromuscular junction, traveling by retrograde axonal transport to the CNS, and disseminating throughout the central nervous system. Clinical features include an incubation of 20--90 days, followed by a prodrome (fever, malaise, paresthesia at the bite site), then acute neurologic phase: furious rabies (80\%)---hydrophobia, aerophobia, agitated behavior, autonomic instability, hypersalivation, and seizures; or paralytic rabies (20\%)---ascending flaccid paralysis. Diagnosis is clinical, confirmed by RT-PCR of saliva, CSF, or skin biopsy. Management is supportive once symptoms develop; prevention through pre-exposure prophylaxis and post-exposure prophylaxis (immediate wound cleansing, HRIG, and rabies vaccine) is nearly 100\% effective if administered before symptom onset.

\subsection{Rubella}
\label{sec:Rubella}
Rubella (German measles) is a mild viral illness caused by rubella virus, most significant for its teratogenic effects, with less than 1000 cases reported annually in the US since 2004, and transmission by respiratory droplets and transplacentally. The condition results from the virus replicating in the respiratory epithelium and spreading hematogenously; in congenital rubella, the virus infects the placenta and fetus during the first trimester, causing vasculitis and direct cellular damage. Clinical features include an incubation of 14--21 days. Postnatal rubella: low-grade fever, headache, coryza, and post-auricular/occipital lymphadenopathy, with a pink, maculopapular rash beginning on the face and spreading downward. Congenital rubella syndrome (CRS) includes the classic triad: sensorineural hearing loss (most common, 60--80\%), congenital heart disease, and cataracts, occurring with infection in the first 16 weeks of pregnancy. Diagnosis is by RT-PCR or IgM serology. Management is supportive. Prognosis is excellent for postnatal rubella; congenital rubella syndrome causes permanent disability.

\subsection{Scarlet Fever}
\label{sec:Scarlet Fever}
Scarlet fever is an acute systemic infection caused by exotoxin-producing strains of Group A $\beta$-hemolytic Streptococcus (GAS, Streptococcus pyogenes), predominantly affecting children aged 5--15 years. The condition results from pyrogenic exotoxins (SPE-A, B, C) inducing delayed-type cutaneous hypersensitivity in individuals lacking protective antitoxin antibodies. Clinical features include fever, pharyngitis, strawberry tongue (initially white coat with red papillae, progressing to a beefy-red raw appearance), circumoral pallor, and a diffuse erythematous maculopapular rash with a 'sandpaper' texture that blanches on pressure and accentuates in skin folds (Pastia lines), followed by desquamation. Diagnosis is confirmed by rapid antigen detection test or throat culture. Management requires 10 days of oral penicillin V or amoxicillin (or macrolides in penicillin-allergic patients) to eradicate GAS and prevent acute rheumatic fever and poststreptococcal glomerulonephritis.

\subsection{Varicella-Zoster (Systemic Aspects)}
\label{sec:Varicella-Zoster Systemic}
Varicella-zoster virus is a DNA herpesvirus that causes two distinct clinical syndromes: primary infection (varicella/chickenpox) and reactivation (herpes zoster/shingles), with varicella affecting nearly all individuals in unvaccinated populations and herpes zoster affecting 30\% of individuals over a lifetime. The condition results from VZV entering through the respiratory mucosa, replicating in the tonsils and regional lymph nodes, and causing a primary viremia that seeds the reticuloendothelial system, with a secondary viremia disseminating the virus to the skin, after which VZV establishes latency in dorsal root ganglia. Clinical features of varicella include fever, malaise, and a vesicular rash (centripetal---trunk, face, proximal extremities), with systemic complications including varicella pneumonia, encephalitis, hepatitis, myocarditis, and glomerulonephritis. Diagnosis is clinical, confirmed by PCR of vesicular fluid. Management includes acyclovir (initiated within 24 hours of rash) and varicella immune globulin for exposed high-risk individuals. Prognosis is excellent for immunocompetent children; prevention through varicella vaccine and shingles vaccine is highly effective.

\subsection{West Nile Virus}
\label{sec:West Nile Virus}
West Nile virus is a mosquito-borne flavivirus, first isolated in Uganda (1937) and introduced to the United States in 1999, transmitted by \textit{Culex} species mosquitoes with birds as the primary amplifying host, with 2000--5000 cases reported annually in the US. The condition results from the virus being inoculated by mosquito bite, replicating in dendritic cells, and spreading to lymph nodes and blood, crossing the blood-brain barrier and infecting neurons, causing apoptosis and perivascular inflammation. Clinical features include an incubation of 2--14 days, with 80\% being asymptomatic, 20\% developing West Nile fever (acute fever, headache, malaise, myalgia, and maculopapular rash), and less than 1\% developing neuroinvasive disease (meningitis, encephalitis, and acute flaccid paralysis). Diagnosis is by WNV-specific IgM in CSF or serum. Management is supportive. Prognosis is excellent for non-neuroinvasive disease; neuroinvasive disease has significant morbidity and mortality.

\subsection{Yellow Fever}
\label{sec:Yellow Fever}
Yellow fever is a mosquito-borne hemorrhagic fever caused by yellow fever virus (YFV), a flavivirus transmitted by \textit{Aedes} and \textit{Haemogogus} mosquitoes, with 200,000 cases and 30,000 deaths annually in tropical Africa and South America. The condition results from the virus infecting dendritic cells, spreading to lymph nodes and liver, with hepatitis resulting from direct viral cytopathology of hepatocytes (midzonal necrosis, Councilman bodies), and the pathogenesis of hemorrhage including hepatic coagulopathy and thrombocytopenia. Clinical features include an incubation of 3--6 days, with three phases: acute (fever, chills, headache, myalgia, conjunctival injection, Faget's sign), remission (hours to 1 day), and intoxication phase (15--25\% progress): recurrence of fever, jaundice, hepatic failure, renal failure, hemorrhagic manifestations (hematemesis---``black vomit''), and shock. Diagnosis is by RT-PCR or antigen detection and serology. Management is supportive. Prognosis: mortality in the intoxication phase is 20--50\%.

\subsection{Zika Virus}
\label{sec:Zika Virus}
Zika virus is a mosquito-borne flavivirus transmitted primarily by \textit{Aedes aegypti} and \textit{A. albopictus}, with other routes including sexual, vertical, blood transfusion, and laboratory exposure, with major outbreaks in Micronesia (2007), French Polynesia (2013--2014), and South America (2015--2016). The condition results from the virus replicating in dendritic cells and spreading hematogenously, with tropism for neural progenitor cells, inducing apoptosis and impairing neurogenesis---the basis for congenital Zika syndrome. Clinical features include an incubation of 3--14 days, with 80\% being asymptomatic; symptomatic disease presents with self-limited febrile illness with maculopapular rash (pruritic, descending), conjunctivitis (non-purulent), arthralgia, and headache. Congenital Zika syndrome includes microcephaly, intracranial calcifications, ventriculomegaly, and arthrogryposis. Guillain-Barré syndrome is associated with Zika infection. Diagnosis is by RT-PCR of serum or urine within 14 days of symptom onset, with serology for confirmation. Management is supportive. Prognosis is excellent for acute illness in adults; congenital infection has significant morbidity.

